\printedpage{24}
% head: A. Grothendieck : Dérivateurs, Ch. XII, Édition des Universités de Montpellier II et Paris VII
% foot: Version 7 décembre 2002
Indeed, if $u'$ is in $x\backslash X$ and $v'$ in $y\backslash Y$, and if $x'$, $y'$ are their images in $X$, $Y$ respectively, then
\[
u'\backslash(x\backslash X)\simeq x'\backslash X,\qquad v'\backslash(y\backslash Y)\simeq y'\backslash Y,
\]
and the arrow between the two is identified with the localization arrow
\[
x'\backslash X\longrightarrow y'\backslash Y.
\]

\textbf{Proposition 5.} \textit{Let $f : X\longrightarrow Y$ be in $\Cat$, suppose $f$ $W$-proper or $W$-smooth, and with $W$-aspheric fibers, then $f\in W$ (and even $f$ is} \textup{universally} \textit{in $W$, i.e. it is in $W$ and remains so under every base change).}

\leavevmode\msp{page 30}As the hypotheses are stable under base change, it suffices to prove $f\in W$. Suppose e.g. $f$ proper, let us prove that $f$ is $W$-\textit{aspheric}, i.e. that the $X/y$ ($y\in\Ob\,Y$) are $W$-aspheric. But $X_y\longrightarrow X/y$, being $W$-coaspheric by the properness hypothesis, is in $W$ (by W(2 bis)), and since $X_y$ is $W$-aspheric by hypothesis, $X/y$ too. But by W(2 bis), $f$ $W$-aspheric $\Longrightarrow f\in W$. The case $f$ smooth is dual.

\textbf{Remark.} Thus, if $f$ is proper, if it has $W$-aspheric fibers it is ``universally $W$-aspheric'', i.e. is and remains $W$-aspheric under every base change. Conversely, if $f$ is universally in $W$ (not necessarily proper), then in particular its fibers are $W$-aspheric. But that of course doesn't imply that $f$ is proper, or smooth (\footnote[44]{But ``universally in $W$'' already implies ``universally $W$-aspheric'', since the notion of $W$-asphericity is defined by \texttt{[the]} proposition 1 a).}).

\textbf{Corollary.} \textit{If $W$ satisfies} \textup{W(2) (not just W(2 bis))}\textit{, then for an $S$-morphism $X\longrightarrow Y$ of categories which are both $W$-proper over $S$, or both $W$-smooth, if $f_s : X_s\longrightarrow Y_s$ is in $W$ for every $s\in S$, then $f\in W$.}

\leavevmode\msp{page 31}Suppose e.g. $f$ proper, let us consider, for $s\in\Ob\,S$
\[
\begin{tikzcd}[row sep=large, column sep=large]
X_s \arrow[r, "f_s"] \arrow[d, "i_s^X"'] & Y_s \arrow[d, "i_s^Y"] \\
X/s \arrow[r, "f/s"] & Y/s,
\end{tikzcd}
\]
as $X$, $Y$ \texttt{[are]} proper, the vertical arrows are in $W$ (much less than properness would suffice \dots), hence $f_s\in W\Longrightarrow f/s\in W$. Thus, ``$f$ is in $W$ locally over $S$'', which implies $f\in W$ by W(2).

\printedpage{25}
% head: A. Grothendieck : Dérivateurs, Ch. XII, Édition des Universités de Montpellier II et Paris VII
% foot: Version 7 décembre 2002
\textbf{Proposition 6.} \textit{Let}
\[
\begin{tikzcd}[row sep=large, column sep=large]
X & X' \arrow[l, "q"'] \arrow[d, "f'"] \\
S \arrow[from=u, "f"'] & S', \arrow[l, "p"]
\end{tikzcd}
\]
\textit{be cartesian in} $\Cat$. \textit{If $f$} \texttt{[is]} \textit{$W$-aspheric, $p$ $W$-smooth, then $f'$ is $W$-aspheric.}

We must prove \texttt{[that]} $X'\times_{S'}S'/s'$ is $W$-aspheric ($s'\in\Ob\,S'$). Let $s=p(s')$, we have
\[
\begin{tikzcd}[column sep=small]
S & & S' \arrow[ll, "p"'] & \\
& S/s \arrow[ul] & & S'/s', \arrow[ul] \arrow[ll, "p_0"]
\end{tikzcd}
\]
with $p_0$ $W$-smooth (lemma p.~28 \texttt{[rather 29]}) and with $W$-aspheric fibers, hence a universal $W$-equivalence (prop.~5). As the inverse image of $X$ over $S/s$ is $W$-aspheric (hypothesis for $f$), it \texttt{[follows]}

\leavevmode\msp{page 32}that in the cartesian square opposite \texttt{[below]}
\[
(*)\qquad\qquad
\begin{tikzcd}[row sep=large, column sep=large]
X/s \arrow[d] & X'/s' \arrow[l, "W"'] \arrow[d] \\
S/s & S'/s' \arrow[l, "W"']
\end{tikzcd}
\]
the horizontal arrows are in $W$, and the first vertical arrow is so by the hypothesis on $f$, hence the second also, hence $X'/s'$ \texttt{[is]} $W$-aspheric.

\textbf{Corollary} (generalizes \texttt{[the]} prop.~6, obtained by taking $Y=S$). \textit{Let $p : S'\longrightarrow S$ be} \texttt{[a]} \textit{$W$-smooth morphism in} $\Cat$. \textit{Then the base change functor}
\[
\Cat/S\longrightarrow\Cat/S'
\]
\textit{is compatible with the localizers $W^{\mathrm{g}}_S$, $W^{\mathrm{g}}_{S'}$, i.e. transforms $f : X\longrightarrow Y$ which is in $W$} locally \textit{over $S$, into $f' : X'\longrightarrow Y'$ in $W$ locally over $S'$.}

This time, we look at the diagram similar to $(*)$
\[
\begin{tikzcd}[column sep=small]
X/s \arrow[dd] \arrow[dr, "f/s"] & & X'/s' \arrow[ll, "W"'] \arrow[dd] \arrow[dr, "f'/s'"] & \\
& Y/s \arrow[dl] & & Y'/s' \arrow[ll, "W"', crossing over] \arrow[dl] \\
S/s & & S'/s', \arrow[ll, "W"'] &
\end{tikzcd}
\]

\printedpage{26}
% head: A. Grothendieck : Dérivateurs, Ch. XII, Édition des Universités de Montpellier II et Paris VII
% foot: Version 7 décembre 2002
where the horizontal arrows are in $W$, hence $f/s\in W$ (hypothesis) implies $f'/s'\in W$.

\leavevmode\msp{page 33}Dually, if $p$ is $W$-proper, the inverse image functor is compatible with the localizers $W^{\mathrm{d}}_S$, $W^{\mathrm{d}}_{S'}$.

If $f : X\longrightarrow S$ is $W$-proper, we are going to define a canonical functor
\[
\begin{aligned}
\hot_*^{X/S,W} : S &\quad\longrightarrow\quad \Hot_W\overset{\text{def}}{=}(\Cat)W^{-1}\\
\text{given on objects by}\qquad s &\quad\longmapsto\quad \hot_W(X_s).
\end{aligned}
\]
We must define, for an arrow $u$ in $S$, $u : s\longrightarrow s'$, an arrow
\[
\hot_W(X_s)\longrightarrow\hot_W(X_{s'})
\]
in $\Hot_W$ - arrow which we can denote by $\hot_*^{X/S,W}(u)$. We define it by the diagram
\[
\begin{tikzcd}[row sep=large, column sep=large]
X_s \arrow[r, dotted] \arrow[d, "i_s^X"] & X_{s'} \arrow[d, "i_{s'}^X"] \\
X/s \arrow[r, "X/u"] & X/s',
\end{tikzcd}
\]
where the vertical arrows are $W$-coaspheric, hence in $W$. They are therefore isomorphisms in $\Hot_W$, whence the unique dotted arrow in $\Hot_W$, making the square commutative. From this definition it is immediate that
\[
u\longmapsto\hot_W(u)
\]
is compatible with composition, and units.

This functor is thus defined assuming

\leavevmode\msp{page 34}only that the
\[
X_y\longrightarrow X/y
\]
are in $W$, which is distinctly weaker than proper. If we suppose that this hypothesis is satisfied ``universally'' (i.e. remains satisfied after every base change $S'\longrightarrow S$) (\footnote[45]{We shall say that $f$ is $W$-\textit{preproper.}}), then we therefore have, for every $S'$ in $\Cat/S$, a functor
\[
\hot_*^{X'/S',W} : S'\longrightarrow\Hot_W,
\]
and we check at once that it is none other than the composite in

\printedpage{27}
% head: A. Grothendieck : Dérivateurs, Ch. XII, Édition des Universités de Montpellier II et Paris VII
% foot: Version 7 décembre 2002
\[
(*)\qquad\qquad
\begin{tikzcd}[column sep=small]
S' \arrow[rr, "p"] \arrow[dr, "\hot_W^{X'/S'}"'] & & S \arrow[dl, "\hot_W^{X/S}"] \\
& \Hot_W &
\end{tikzcd}
\]
\texttt{[rather} $\hot_*^{X'/S',W}$ \texttt{resp.} $\hot_*^{X/S,W}$\texttt{]}. We have dual observations if $X$ is smooth over $S$, more generally if the inclusions
\[
X_s\longrightarrow s\backslash X
\]
are in $W$. We then find a contrafunctor
\[
\left\{
\begin{aligned}
S^\circ &\quad\longrightarrow\quad \Hot_W\\
s &\quad\longmapsto\quad \hot_W(X_s)\\
u &\quad\longmapsto\quad \hot^*_{X/S,W}(u)\ \text{or}\ \hot^*_{X/S}(u)
\end{aligned}
\right.
\]

\leavevmode\msp{page 35}and a compatibility similar to $(*)$, if the hypothesis is satisfied \textit{universally} over $S$, e.g. $f$ $W$-smooth.

\textbf{Proposition 7.} \textit{Suppose $f : X\longrightarrow S$ both $W$-proper and $W$-smooth, more generally that the morphisms}
\[
X_s\longrightarrow X/s,\qquad X_s\longrightarrow s\backslash X\qquad(s\in S)
\]
\textit{are in $W$, and that this hypothesis remains satisfied after every base change. Then the corresponding functors} % [sic] verifiée -> vérifiée, accent missing
\[
\begin{aligned}
\lambda &\;=\;\hot_*^{X/S,W} : &&\quad S \quad\longrightarrow\quad \Hot_W\\
\lambda^\circ &\;=\;\hot^*_{X/S,W} : &&\quad S^\circ \quad\longrightarrow\quad \Hot_W
\end{aligned}
\]
\textit{transform every arrow $u$ of $S$ into} \texttt{[an]} \textit{isomorphism of} $\Hot_W$\textit{, and we have}
\[
\lambda^\circ(u)=\lambda(u)^{-1}\qquad\textit{dans}\quad\Hot_W,
\]
\textit{i.e. $\lambda(u)$ and $\lambda^c irc(u)$ are inverses of each other.} % [sic] λ^c irc(u) -> λ°(u), command \circ misprinted as superscript c followed by "irc"

\textsc{Proof}. The datum of $u : s_0\longrightarrow s_1$ is equivalent to that of
\[
S'=\mathbf{\Delta}^1\longrightarrow S,
\]
whence
\[
X'=X\times_S S'\qquad\text{over } S'=\mathbf{\Delta}^1,
\]
with the two fibers
\[
X'_0\simeq X_{s_0},\qquad X'_1\simeq X_{s_1}.
\]
The hypothesis ensuring the existence of $\lambda$ is expressed by the condition that
\[
X'_1\longrightarrow X'/1\simeq X'
\]
\printedpage{28}
% head: A. Grothendieck : Dérivateurs, Ch. XII, Édition des Universités de Montpellier II et Paris VII
% foot: Version 7 décembre 2002
\msp{page 36}is in $W$, so that the diagram
\[
X'_0\xrightarrow{i_0}X'\xleftarrow[W]{i_1}X'_1
\]
defines an arrow $\lambda(u) : X'_0\longrightarrow X'_1$ in $\Hot_W$, namely $\lambda(u)=i_1^{-1}i_0$. The hypothesis which ensures the existence of $\lambda^\circ$ is expressed by the symmetric condition $i_0\in W$, which makes it possible to define the arrow $\lambda^\circ(u) : X'_1\longrightarrow X'_0$ by $\lambda^\circ(u)=i_0^{-1}i_1$. When the two conditions are satisfied simultaneously, we therefore see that $\lambda(u)$, $\lambda^\circ(u)$ are inverses of each other.

\textbf{Remarks.} To define $\lambda$ ``universally'', it was not indispensable to consider the $X/s$ and the inclusions $X_s\longrightarrow X/s$, instead of that it sufficed to look at the base changes $S'=\mathbf{\Delta}^1\xrightarrow{u}S$, and to make the hypothesis that for every such base change, the inclusion $X'_1\xrightarrow{i_1}X'=X_{(\mathbf{\Delta}^1,u)}$ is in $W$. This hypothesis is automatically stable under arbitrary base change. (The $W$-properness of $f : X\longrightarrow S$ means that $X'_1\longrightarrow X'$ is not only in $W$, but even $W$-coaspheric.) Dually, for the ``universal'' definition of $\lambda^\circ$, it would suffice that for every base change $\mathbf{\Delta}^1\longrightarrow S$ as above, the inclusion $X'_0\xrightarrow{i_0}X'$ is in $W$.

\leavevmode\msp{page 37}The weakest reasonable notion of $W$-fibered, seems to me to be the one where we suppose both things at once. The hypotheses we have made up to now, excluding the strong form of $W$-properness or of $W$-smoothness, always amount to supposing that the \textit{base change by $f : X\longrightarrow S$} transforms certain arrows in $\Cat/S$ into \texttt{[an element of]} $W$ : To define $\lambda$, these arrows are
\[
\text{either } \{s\}\longrightarrow S/s,\qquad \text{or } \{1\}\longrightarrow\mathbf{\Delta}^1(\longrightarrow S),
\]
to define $\lambda^\circ$, they are
\[
\text{either } \{s\}\longrightarrow s\backslash S,\qquad \text{or } \{0\}\longrightarrow\mathbf{\Delta}^1(\longrightarrow S),
\]
the ``either'' depending on the definition we have adopted. In these examples, it is always a matter of the inclusion of the final object of a category $S'$, into the said category, or conversely of the initial object. This shows moreover (if we require universality), that the first hypothesis (ensuring $\lambda$) is \textit{exactly} that (case of the inclusion of a final object $s'$ - since for such an object, we have $S'/s'\simeq S'$), whereas the second hypothesis requires less, only for the case of $\mathbf{\Delta}^1$. And dually for $\lambda'$ \texttt{[rather $\lambda^\circ$]}. But it should be noted that in practice, the most convenient way, and by far,

\leavevmode\msp{page 38}to verify such properties, is either by the $W$-properness of $f$, or by $W$-smoothness - which reduces to verifying that certain well-specified categories are $W$-aspheric, without having to worry about base changes and about verifying that the arrows constructed in terms of arbitrary base changes are in $W$. And the use of $W$-properness and $W$-smoothness is all the more congenial, as these notions are stable under composition,